\noindent\underline{Proof of Proposition 3.3.}
We may restrict to the case $X\ne\mathbb P^r$.  The two properties
asserted to be equivalent are independent of the particular choice of
coordinates, so we may describe them using coordinates chosen in a
very special way.  Since $X$ is a smooth subvariety of $\mathbb P^r$
of dimension $n\geq1$, choose homogeneous coordinates
$X_0,\ldots,X_r$ so that:
\begin{enumerate}
\item[3.3.1.] the point $x_0\in\mathbb P^r$ is $(1,0,\ldots,0)$;
\item[3.3.2.] putting $x_i=X_i/X_0$ for $i=1,\ldots,r$, so that
$\widehat{\mathcal O}_{\mathbb P^r,x_0}$ identifies with
$k[[x_1,\ldots,x_r]]$, the ideal defining the trace of $X$ in
$\operatorname{Spec}\widehat{\mathcal O}_{\mathbb P^r,x_0}$ is
generated by elements $g_i$ $(i=1,\ldots,r-n)$ which, in the
completion, have the form
\[
g_i=x_{n+i}-h_i(x_1,\ldots,x_n),
\qquad i=1,\ldots,r-n,
\]
where the $h_i$ are power series in $x_1,\ldots,x_n$ alone and lie in
the square of the maximal ideal of $k[[x_1,\ldots,x_n]]$;
\item[3.3.3.] under (3.1.1), the hyperplane $H$ comes from
$g_{r-n}$.
\end{enumerate}

Put $V=\operatorname{Spec}(\widehat{\mathcal O}_{X,x_0})$.  It follows
from 3.3.2 that there is a trivialization
\[
(I/I^2)|_V\simeq\mathcal O_V^{r-n}
\]
given by
\begin{equation}\tag{3.3.4}
(\Lambda_1,\ldots,\Lambda_{r-n})\longmapsto
\sum_{i=1}^{r-n}\Lambda_i g_i.
\end{equation}
Together with 3.3.3 this gives a neighborhood $W$ of $(x_0,H)$ in
$\mathbb P(N)|_V$ and an isomorphism
$V\times\mathbb A^{r-n-1}\simeq W$, given by
\begin{equation}\tag{3.3.5}
\begin{aligned}
V\times\mathbb A^{r-n-1}&\longrightarrow\mathbb P(I/I^2)|_V,\\
(v,(\lambda_1,\ldots,\lambda_{r-n-1}))&\longmapsto
\left(v,\left[g_{r-n}+\sum_{i=1}^{r-n-1}\lambda_i g_i\right]\right),
\end{aligned}
\end{equation}
where $\lambda_i=\Lambda_i/\Lambda_{r-n}$.

In particular, the completion of the local ring of $\mathbb P(N)$ at
$(x_0,H)$ identifies with
\[
k[[x_1,\ldots,x_n,\lambda_1,\ldots,\lambda_{r-n-1}]].
\]
We now express $\varphi$ in these coordinates at the level of
completions:
\begin{equation}\tag{3.3.6}
\varphi(x,\lambda)
=\text{the hyperplane with equation }
\sum_{i=0}^{r}\varphi_i(x,\lambda)T_i=0.
\end{equation}
Since this image hyperplane contains the point $x$, whose homogeneous
coordinates are
\begin{equation}\tag{3.3.7}
x=(1,x_1,\ldots,x_n,h_1(x),\ldots,h_{r-n}(x)),
\end{equation}
we have
\begin{equation}\tag{3.3.8}
\varphi_0(x,\lambda)
=-\sum_{i=1}^{n}x_i\varphi_i(x,\lambda)
 -\sum_{j=1}^{r-n}h_j(x)\varphi_{n+j}(x,\lambda)
\end{equation}
in $\widehat{\mathcal O}_{\mathbb P(N),(x_0,H)}$.  Moreover, by
(3.1.1), if
\[
D_i=\frac{\partial}{\partial x_i}\in
\operatorname{Der}_k(
\widehat{\mathcal O}_{\mathbb P^r,x_0},
\widehat{\mathcal O}_{\mathbb P^r,x_0}),
\]
then for $1\leq i\leq r$,
\begin{equation}\tag{3.3.9}
\varphi_i(x,\lambda)
=D_i\!\left(g_{r-n}+
\sum_{j=1}^{r-n-1}\lambda_jg_j\right)(x)
\quad\text{in }\widehat{\mathcal O}_{\mathbb P(N),(x_0,H)}.
\end{equation}

Recalling that $g_i=x_{n+i}-h_i(x_1,\ldots,x_n)$, one obtains
\begin{equation}\tag{3.3.10}
\varphi_i(x,\lambda)
=-\frac{\partial}{\partial x_i}
\left(h_{r-n}+\sum_{j=1}^{r-n-1}\lambda_jh_j\right),
\qquad 1\leq i\leq n,
\end{equation}
\begin{equation}\tag{3.3.11}
\varphi_{n+j}(x,\lambda)=\lambda_j,
\qquad 1\leq j\leq r-n-1,
\end{equation}
and
\begin{equation}\tag{3.3.12}
\varphi_r(x,\lambda)=1.
\end{equation}
\footnote{The print gives incompatible ranges in (3.3.11).  The
coordinates in (3.3.5), the coefficient formula (3.3.9), and the
identity block of the printed Jacobian matrix force the indices shown
here.}

Let $\mathfrak m$ be the maximal ideal of
$\widehat{\mathcal O}_{\mathbb P(N),(x_0,H)}$.  Since
$h_i\equiv0\pmod{\mathfrak m^2}$ for $i=1,\ldots,r-n$, we find
\begin{equation}\tag{3.3.13}
\varphi_0(x,\lambda)\equiv0\pmod{\mathfrak m^2}
\end{equation}
and
\begin{equation}\tag{3.3.14}
\varphi_i(x,\lambda)
\equiv-\frac{\partial h_{r-n}}{\partial x_i}
\pmod{\mathfrak m^2},
\qquad i=1,\ldots,n.
\end{equation}
